% Stage 2 mathematical core.
% This file is designed to be included after amsmath, amssymb, amsthm,
% mathtools, booktabs, and enumitem have been loaded.

\section{Finite-horizon reliability before sampler selection}
\label{sec:stage2-mathematical-core}

Throughout this section, \(h\) is fixed and \(\pi_h\) denotes the declared
discrete posterior.  The candidate library \(\mathcal A\), the estimands
\(f_1,\ldots,f_R\), the initialization laws \(\nu_a\), and all numerical
thresholds are frozen before the streams considered below are generated.  A
statement about a stationary long-run variance will always say explicitly
that the chain is stationary.  Every other statement is finite-horizon and
retains its initialization law.

\subsection{Finite-horizon risk objects}
\label{subsec:stage2-risk}

For candidate \(a\), let \(P_a\) be a Markov kernel with invariant measure
\(\pi_h\), let \(X_0\sim\nu_a\), and define
\begin{equation}
 \widehat\mu_{a,r,n}
 =
 \frac1n\sum_{k=0}^{n-1} f_r(X_k),
 \qquad
 \mu_r=\pi_h(f_r).
 \label{eq:s2-estimator}
\end{equation}
Whenever \(f_r(X_k)\) is square integrable under the finite path law, set
\begin{align}
 D_{a,r}(n)
 &=
 n\,\operatorname{Var}_{\nu_a}(\widehat\mu_{a,r,n}),
 \label{eq:s2-dispersion}\\
 R_{a,r}(n)
 &=
 n\,\mathbb E_{\nu_a}
 \left[
  \{\widehat\mu_{a,r,n}-\mu_r\}^2
 \right],
 \label{eq:s2-risk}\\
 B_{a,r}(n)
 &=
 \mathbb E_{\nu_a}[\widehat\mu_{a,r,n}]-\mu_r.
 \label{eq:s2-bias}
\end{align}

\begin{proposition}[Finite-horizon bias--variance identity]
\label{prop:s2-bias-variance}
For every \(a,r,n\) for which the preceding quantities are finite,
\begin{equation}
 R_{a,r}(n)
 =
 D_{a,r}(n)+nB_{a,r}(n)^2.
 \label{eq:s2-bias-variance}
\end{equation}
No stationarity or unbiasedness assumption is required.
\end{proposition}

\begin{proof}
Put \(Y=\widehat\mu_{a,r,n}\).  The elementary identity
\[
 \mathbb E_{\nu_a}\{(Y-\mu_r)^2\}
 =
 \operatorname{Var}_{\nu_a}(Y)
 +
 \{\mathbb E_{\nu_a}Y-\mu_r\}^2
\]
follows by writing \(Y-\mu_r=(Y-\mathbb E_{\nu_a}Y)
+(\mathbb E_{\nu_a}Y-\mu_r)\) and observing that the cross term has mean zero.
Multiplication by \(n\) gives \eqref{eq:s2-bias-variance}.
\end{proof}

\begin{proposition}[Two-estimator disagreement]
\label{prop:s2-disagreement}
Let \(U,V\in L^2\) estimate a common scalar target \(\mu\).  Then
\begin{equation}
 \max\{\mathbb E(U-\mu)^2,\mathbb E(V-\mu)^2\}
 \geq
 \frac14\{\mathbb E U-\mathbb E V\}^2.
 \label{eq:s2-disagreement}
\end{equation}
If \(U\) and \(V\) are horizon-\(n\) estimators, the same statement multiplied
by \(n\) holds for their scaled finite-horizon risks.
\end{proposition}

\begin{proof}
Let \(b_U=\mathbb E U-\mu\) and \(b_V=\mathbb E V-\mu\).  By the triangle
inequality and the fact that squared bias does not exceed mean-squared error,
\[
 |\mathbb E U-\mathbb E V|
 \leq |b_U|+|b_V|
 \leq
 2\sqrt{\max\{\mathbb E(U-\mu)^2,\mathbb E(V-\mu)^2\}}.
\]
Squaring proves the claim.
\end{proof}

\subsection{Finite-path non-identifiability under metastability}
\label{subsec:stage2-nonidentifiability}

For \(p\in(0,1)\) and \(\varepsilon\in(0,1]\), consider
\begin{equation}
 P_{p,\varepsilon}
 =
 \begin{pmatrix}
  1-\varepsilon p & \varepsilon p\\
  \varepsilon(1-p) & 1-\varepsilon(1-p)
 \end{pmatrix}.
 \label{eq:s2-two-state}
\end{equation}
Its invariant distribution is
\(\pi_p=(1-p,p)\), and detailed balance follows from
\((1-p)\varepsilon p=p\varepsilon(1-p)\).

\begin{theorem}[Exact finite-path false precision]
\label{thm:s2-false-precision}
Run \(L\) independent copies of \eqref{eq:s2-two-state}, all initialized at
zero, and observe each for \(n\) transitions.  Let
\[
 \mathcal N_{L,n}
 =
 \{\text{none of the \(L\) chains visits state \(1\) during the \(n\)
 transitions}\}.
\]
Then
\begin{equation}
 \mathbb P_{p,\varepsilon}(\mathcal N_{L,n})
 =
 (1-\varepsilon p)^{Ln}.
 \label{eq:s2-no-visit}
\end{equation}
On \(\mathcal N_{L,n}\), every replicate estimate of
\(\pi_p(\{1\})=p\) is zero.  Hence the empirical variance across replicates is
exactly zero, while every replicate has absolute error \(p\) and squared error
\(p^2\).
\end{theorem}

\begin{proof}
A chain that is currently at zero remains there with probability
\(1-\varepsilon p\).  Starting from zero, it avoids state one for all \(n\)
transitions with probability \((1-\varepsilon p)^n\).  Independence of the
\(L\) paths gives \eqref{eq:s2-no-visit}.  On this event all recorded states
are zero, so every empirical mean of the indicator \(f(x)=\mathbf 1\{x=1\}\)
equals zero.  The remaining conclusions are immediate.
\end{proof}

\begin{corollary}[False precision with arbitrarily high probability]
\label{cor:s2-arbitrary-false-precision}
Fix \(p\in(0,1)\), \(L,n\geq1\), and \(\alpha\in(0,1)\).  If
\begin{equation}
 0<\varepsilon
 \leq
 \min\left\{
  1,\,
  \frac{1-(1-\alpha)^{1/(Ln)}}{p}
 \right\},
 \label{eq:s2-epsilon-choice}
\end{equation}
then
\(\mathbb P_{p,\varepsilon}(\mathcal N_{L,n})\geq1-\alpha\).
Thus zero empirical dispersion can occur with probability as close to one as
desired while the error remains exactly \(p\).
\end{corollary}

\begin{proof}
The choice \eqref{eq:s2-epsilon-choice} implies
\(1-\varepsilon p\geq(1-\alpha)^{1/(Ln)}\).  Apply
\eqref{eq:s2-no-visit}.
\end{proof}

For the information-theoretic statement, let
\(\mathbb Q_{p,\varepsilon}^{L,n}\) denote the joint law of the \(L\) observed
paths, including the deterministic initial zero.  We use
\(\|\cdot\|_{\mathrm{TV}}=\sup_A|\mathbb Q(A)-\mathbb Q'(A)|\).

\begin{theorem}[Finite-path indistinguishability]
\label{thm:s2-indistinguishability}
Let \(0<p_0<p_1<1\), let \(\Delta=p_1-p_0\), and use the same
\(\varepsilon\in(0,1]\) for both chains.  Then
\begin{equation}
 \left\|
  \mathbb Q_{p_0,\varepsilon}^{L,n}
  -
  \mathbb Q_{p_1,\varepsilon}^{L,n}
 \right\|_{\mathrm{TV}}
 \leq
 1-(1-\varepsilon p_1)^{Ln}
 \leq
 Ln\varepsilon p_1.
 \label{eq:s2-path-tv}
\end{equation}
Consequently, for every path-measurable estimator \(T\) of the target event
probability,
\begin{equation}
 \max_{i\in\{0,1\}}
 \mathbb Q_{p_i,\varepsilon}^{L,n}
 \left\{
  |T-p_i|\geq\frac{\Delta}{2}
 \right\}
 \geq
 \frac12
 \left(
  1-
  \left\|
   \mathbb Q_{p_0,\varepsilon}^{L,n}
   -
   \mathbb Q_{p_1,\varepsilon}^{L,n}
  \right\|_{\mathrm{TV}}
 \right).
 \label{eq:s2-le-cam}
\end{equation}
In particular, take \(p_0=1/4\), \(p_1=3/4\), and any
\(\eta\in(0,1)\).  If
\[
 \varepsilon
 \leq
 \min\left\{1,\frac{\eta}{Ln p_1}\right\},
\]
then the two length-\(n\) experiments have total variation distance at most
\(\eta\), although their target event probabilities differ by \(1/2\), and
the right side of \eqref{eq:s2-le-cam} is at least
\((1-\eta)/2\).
\end{theorem}

\begin{proof}
Let \(x_\circ\) be the joint observation in which every recorded state of every
chain is zero.  Under parameter \(p_i\),
\[
 \mathbb Q_{p_i,\varepsilon}^{L,n}(\{x_\circ\})
 =
 (1-\varepsilon p_i)^{Ln}.
\]
For probability measures on a countable sample space,
\(\|Q-Q'\|_{\mathrm{TV}}=1-\sum_x\min\{Q(x),Q'(x)\}\).  Retaining only the
common atom \(x_\circ\) in this sum and using \(p_0<p_1\) gives
\[
 \| \mathbb Q_{p_0,\varepsilon}^{L,n}
      -\mathbb Q_{p_1,\varepsilon}^{L,n}\|_{\mathrm{TV}}
 \leq
 1-(1-\varepsilon p_1)^{Ln}.
\]
Bernoulli's inequality \(1-(1-x)^m\leq mx\), for \(x\in[0,1]\), proves the
second bound in \eqref{eq:s2-path-tv}.

To prove \eqref{eq:s2-le-cam}, form the test
\(\phi=\mathbf 1\{T\geq(p_0+p_1)/2\}\).  If
\(|T-p_0|<\Delta/2\), then \(\phi=0\), and if
\(|T-p_1|<\Delta/2\), then \(\phi=1\).  Hence
\[
 \mathbb Q_{p_0,\varepsilon}^{L,n}
 \{|T-p_0|\geq\Delta/2\}
 +
 \mathbb Q_{p_1,\varepsilon}^{L,n}
 \{|T-p_1|\geq\Delta/2\}
 \geq
 \mathbb Q_{p_0,\varepsilon}^{L,n}(\phi=1)
 +
 \mathbb Q_{p_1,\varepsilon}^{L,n}(\phi=0).
\]
The sum of the two testing errors is at least one minus the total variation
distance.  At least one estimation error probability is therefore at least
half that quantity.  The numerical specialization follows from
\eqref{eq:s2-path-tv}.
\end{proof}

\begin{corollary}[No uniformly reliable nontrivial finite-path certificate]
\label{cor:s2-no-certificate}
In the setting of Theorem~\ref{thm:s2-indistinguishability}, fix
\(\tau<\Delta/2\).  Let \(C\) be the event that a path-measurable diagnostic
certifies an estimator \(T\).  Suppose, for both \(i=0,1\),
\begin{align}
 \mathbb Q_{p_i,\varepsilon}^{L,n}(C)
 &\geq 1-\beta,
 \label{eq:s2-cert-power}\\
 \mathbb Q_{p_i,\varepsilon}^{L,n}
 \{C\cap(|T-p_i|>\tau)\}
 &\leq\alpha.
 \label{eq:s2-cert-error}
\end{align}
Then
\begin{equation}
 \alpha+\beta
 \geq
 \frac12
 \left(
  1-
  \left\|
   \mathbb Q_{p_0,\varepsilon}^{L,n}
   -
   \mathbb Q_{p_1,\varepsilon}^{L,n}
  \right\|_{\mathrm{TV}}
 \right).
 \label{eq:s2-cert-impossibility}
\end{equation}
Thus, when the path laws are close, a uniformly low false-certification
probability forces a substantial abstention probability.  Acceptance,
effective sample size, replicate dispersion, or any other statistic measurable
from the same finite paths cannot evade this conclusion without additional
structural assumptions.
\end{corollary}

\begin{proof}
For each \(i\),
\[
 \mathbb Q_{p_i,\varepsilon}^{L,n}(|T-p_i|>\tau)
 \leq
 \mathbb Q_{p_i,\varepsilon}^{L,n}(C^c)
 +
 \mathbb Q_{p_i,\varepsilon}^{L,n}
 \{C\cap(|T-p_i|>\tau)\}
 \leq\alpha+\beta.
\]
Apply the Le Cam bound \eqref{eq:s2-le-cam}; replacing
\(\Delta/2\) by the smaller tolerance \(\tau\) can only increase the error
event.
\end{proof}

\subsection{Exact rare-mode proxy failure}
\label{subsec:stage2-proxy}

Let \(\mathsf M_A,\mathsf M_B\) be disjoint, let \(\pi_A,\pi_B\) be
probability measures supported on the respective sets, and define
\begin{equation}
 \pi_p=(1-p)\pi_A+p\pi_B,
 \qquad
 q_w=(1-w)\pi_A+w\pi_B,
 \qquad
 0<w<p<1.
 \label{eq:s2-mode-mixtures}
\end{equation}

\begin{theorem}[Rare-mode proxy formulas]
\label{thm:s2-proxy}
Let \(P_w\) be the independence Metropolis--Hastings kernel with target
\(\pi_p\), proposal \(q_w\), and stationary initialization.  Then
\begin{align}
 \operatorname{Acc}(p,w)
 &=
 1-p+w,
 \label{eq:s2-acceptance}\\
 \sigma_{P_w}^2(\mathbf 1_{\mathsf M_B})
 &=
 p(1-p)\left(\frac{2p}{w}-1\right),
 \label{eq:s2-proxy-lrv}\\
 D_{\mathrm{KL}}(q_w\Vert\pi_p)
 &=
 (1-w)\log\frac{1-w}{1-p}
 +
 w\log\frac{w}{p}.
 \label{eq:s2-reverse-kl}
\end{align}
\end{theorem}

\begin{proof}
The importance ratio \(\pi_p/q_w\) is \((1-p)/(1-w)\) on
\(\mathsf M_A\) and \(p/w\) on \(\mathsf M_B\).  Because \(w<p\), proposals
within a mode and proposals from \(\mathsf M_A\) to \(\mathsf M_B\) are
accepted.  A proposal from \(\mathsf M_B\) to \(\mathsf M_A\) is accepted with
probability \(w(1-p)/\{p(1-w)\}\).  The induced two-state mode chain therefore
has
\[
 P_w(A,B)=w,
 \qquad
 P_w(B,A)=\frac{w(1-p)}p,
\]
stationary probability \(p\) for \(B\), and nonconstant eigenvalue
\(\lambda_w=1-w/p\).  Conditional acceptance is one in mode \(A\) and
\(w/p\) in mode \(B\), proving \eqref{eq:s2-acceptance}.

For \(g=\mathbf 1_{\mathsf M_B}-p\),
\[
 \operatorname{Cov}_{\pi_p}\{g(X_0),g(X_k)\}
 =
 p(1-p)\lambda_w^k.
\]
Since \(0<\lambda_w<1\), summing the stationary covariance series gives
\[
 \sigma_{P_w}^2(\mathbf 1_{\mathsf M_B})
 =
 p(1-p)
 \left\{1+2\sum_{k=1}^{\infty}\lambda_w^k\right\}
 =
 p(1-p)\frac{1+\lambda_w}{1-\lambda_w},
\]
which simplifies to \eqref{eq:s2-proxy-lrv}.  Finally,
\eqref{eq:s2-reverse-kl} is obtained by integrating
\(\log(dq_w/d\pi_p)\) under \(q_w\).
\end{proof}

\begin{corollary}[Precise favourable-proxy limits]
\label{cor:s2-proxy-limits}
Let \(p_j\downarrow0\) and \(0<w_j<p_j\).
\begin{enumerate}[label=\textnormal{(\roman*)}]
\item If \(w_j/p_j\to0\), then
\[
 \operatorname{Acc}(p_j,w_j)\to1,
 \qquad
 D_{\mathrm{KL}}(q_{w_j}\Vert\pi_{p_j})\to0,
 \qquad
 \frac{\sigma_{P_{w_j}}^2(\mathbf 1_{\mathsf M_B})}
      {p_j(1-p_j)}
 =
 \frac{2p_j}{w_j}-1
 \to\infty.
\]
\item If, in addition, \(w_j/p_j^2\to0\), then the absolute LRV diverges:
\[
 \sigma_{P_{w_j}}^2(\mathbf 1_{\mathsf M_B})\to\infty.
\]
For example, \(w_j=p_j^3\) satisfies all of these conclusions.
\end{enumerate}
For fixed \(p>0\), in contrast, \(w\downarrow0\) gives acceptance
\(1-p\) and reverse KL \(-\log(1-p)\), not one and zero.  Exact reverse-KL
minimization over \(w\in[0,1]\) still has the correct minimizer \(w=p\).
\end{corollary}

\begin{proof}
The acceptance and relative-LRV limits follow directly from
\eqref{eq:s2-acceptance} and \eqref{eq:s2-proxy-lrv}.  For reverse KL, the
first term in \eqref{eq:s2-reverse-kl} converges to zero.  The second also
converges to zero: because \(0<w_j<p_j<1\),
\[
 0\leq -w_j\log(w_j/p_j)
 \leq -w_j\log w_j\longrightarrow0.
\]
Nonnegativity of KL then yields the claimed limit.  Finally,
\[
 \sigma_{P_{w_j}}^2(\mathbf 1_{\mathsf M_B})
 =
 2(1-p_j)\frac{p_j^2}{w_j}-p_j(1-p_j),
\]
which diverges when \(w_j/p_j^2\to0\).  The final statements follow by direct
calculation and strict convexity of the Bernoulli KL divergence.
\end{proof}
\subsection{A statistically calibrated exploration gate}
\label{subsec:stage2-exploration}

Let \(\mathcal F_{\mathrm{freeze}}\) contain the target-validity, pilot, mode,
estimand, candidate-library, and initialization information.  Conditional on
\(\mathcal F_{\mathrm{freeze}}\), all these objects are fixed.  For each
tested candidate--estimand pair \((a,r)\), define
\begin{equation}
 z_{a,r}(n)
 =
 \mathbb P_{\nu_a}
 \{N_{a,r}^{\mathrm{cross}}(n)=0
   \mid\mathcal F_{\mathrm{freeze}}\}.
 \label{eq:s2-zero-cross-prob}
\end{equation}
The \(L_{\mathrm{adq}}\) adequacy paths for a fixed candidate are independent
replicates of length \(n\) with initialization \(\nu_a\).  Dependence between
different candidate--estimand pairs is allowed.

For \(0\leq k\leq L\) and \(\eta\in(0,1)\), define the one-sided
Clopper--Pearson upper bound
\begin{equation}
 U_{\mathrm{CP}}(k,L;\eta)
 =
 \begin{cases}
  1, & k=L,\\[2mm]
  F_{\mathrm{Beta}(k+1,L-k)}^{-1}(1-\eta), & 0\leq k<L.
 \end{cases}
 \label{eq:s2-cp}
\end{equation}
Equivalently, for \(k<L\), \(U_{\mathrm{CP}}\) is the unique \(u\in(0,1)\)
satisfying
\[
 \sum_{j=0}^{k}\binom{L}{j}u^j(1-u)^{L-j}=\eta.
\]

\begin{theorem}[Simultaneous Gate-2 calibration]
\label{thm:s2-gate2}
Let \(\mathcal I_{\mathrm{adq}}\) be the set of tested
candidate--mandatory-nonsmooth-estimand pairs and put
\(M_{\mathrm{adq}}=|\mathcal I_{\mathrm{adq}}|\).  Allocate numbers
\(\eta_{a,r}>0\) satisfying
\[
 \sum_{(a,r)\in\mathcal I_{\mathrm{adq}}}\eta_{a,r}
 \leq\delta_{\mathrm{adq}}.
\]
If \(K_{a,r}\) of the \(L_{\mathrm{adq}}\) paths have zero crossings, then,
conditionally on \(\mathcal F_{\mathrm{freeze}}\),
\[
 K_{a,r}\sim
 \operatorname{Binomial}(L_{\mathrm{adq}},z_{a,r}(n)).
\]
Moreover,
\begin{equation}
 \mathbb P\left\{
  z_{a,r}(n)
  \leq
  U_{a,r}:=
  U_{\mathrm{CP}}(K_{a,r},L_{\mathrm{adq}};\eta_{a,r})
  \ \text{for all }(a,r)\in\mathcal I_{\mathrm{adq}}
  \,\middle|\,\mathcal F_{\mathrm{freeze}}
 \right\}
 \geq1-\delta_{\mathrm{adq}}.
 \label{eq:s2-cp-simultaneous}
\end{equation}
Consequently, a candidate declared eligible only when
\begin{equation}
 U_{a,r}\leq z_r^{\max}
 \quad
 \text{for every mandatory nonsmooth \(r\)}
 \label{eq:s2-gate2-pass}
\end{equation}
satisfies \(z_{a,r}(n)\leq z_r^{\max}\) for all of those estimands on the
simultaneous event \eqref{eq:s2-cp-simultaneous}.
\end{theorem}

\begin{proof}
For a fixed pair, the zero-crossing indicators are conditionally independent
Bernoulli variables with success probability \(z_{a,r}(n)\), proving the
binomial statement.  The defining inversion of the exact binomial test gives
\[
 \mathbb P\{z_{a,r}(n)>U_{a,r}
 \mid\mathcal F_{\mathrm{freeze}}\}
 \leq\eta_{a,r}.
\]
The union bound over \(\mathcal I_{\mathrm{adq}}\) proves
\eqref{eq:s2-cp-simultaneous}.  On that event,
\eqref{eq:s2-gate2-pass} implies
\(z_{a,r}(n)\leq U_{a,r}\leq z_r^{\max}\).
\end{proof}

\begin{corollary}[Zero observed failures and required replicate count]
\label{cor:s2-gate2-count}
If \(K_{a,r}=0\), then
\begin{equation}
 U_{\mathrm{CP}}(0,L;\eta)
 =
 1-\eta^{1/L}.
 \label{eq:s2-cp-zero}
\end{equation}
Thus zero observed zero-crossing paths can establish
\(z_{a,r}(n)\leq z^{\max}\) only if
\begin{equation}
 L
 \geq
 \left\lceil
  \frac{\log\eta}{\log(1-z^{\max})}
 \right\rceil.
 \label{eq:s2-required-adq}
\end{equation}
In particular, with \(L=12\) and unadjusted one-sided confidence \(0.95\),
\[
 U_{\mathrm{CP}}(0,12;0.05)
 =
 1-0.05^{1/12}
 =
 0.220922\ldots.
\]
Hence \(0/12\) does not establish \(z\leq0.10\).
\end{corollary}

\begin{proof}
For \(k=0\), the binomial lower tail is \((1-u)^L\).  Setting it equal to
\(\eta\) gives \eqref{eq:s2-cp-zero}.  Solving
\(1-\eta^{1/L}\leq z^{\max}\) for \(L\) gives
\eqref{eq:s2-required-adq}; both logarithms are negative.
\end{proof}

\begin{remark}[Entry and exit counts are severity measures]
\label{rem:s2-entry-exit}
The eligibility rule \eqref{eq:s2-gate2-pass} uses only the simultaneous
zero-crossing-probability bounds.  Mean entry and exit totals are reported as
secondary severity measures and are not compared with a point threshold of
five.  If a later protocol makes a replicate-level count \(N\) an eligibility
condition, it must use a simultaneous lower confidence bound.  For example,
if \(0\leq N\leq c_n\) almost surely over independent paths, Hoeffding's
inequality gives, simultaneously over \(M\) declared count summaries,
\[
 \mathbb E N
 \geq
 \overline N
 -
 c_n\sqrt{\frac{\log(M/\delta)}{2L}}
\]
with probability at least \(1-\delta\).  This conservative bound is displayed
before any count threshold is adopted.
\end{remark}

\subsection{Horizon adequacy and qualification of LRV language}
\label{subsec:stage2-horizon}

For a stationary chain and a centred square-integrable \(f\), write
\(\gamma_k(f)=\operatorname{Cov}_{\pi_h}\{f(X_0),f(X_k)\}\).  Define
\[
 D(n)
 =
 n\,\operatorname{Var}_{\pi_h}
 \left(\frac1n\sum_{j=0}^{n-1}f(X_j)\right).
\]

\begin{theorem}[A quantitative stationary plateau envelope]
\label{thm:s2-lrv-envelope}
Assume that the chain is stationary and
\begin{equation}
 |\gamma_k(f)|\leq C\rho^k,
 \qquad k\geq1,
 \qquad C<\infty,\quad0<\rho<1.
 \label{eq:s2-cov-envelope}
\end{equation}
Then the stationary long-run variance
\[
 \sigma_P^2(f)
 =
 \gamma_0(f)+2\sum_{k=1}^{\infty}\gamma_k(f)
\]
exists absolutely, and
\begin{equation}
 |D(n)-\sigma_P^2(f)|
 \leq
 2C
 \left[
  \frac{\rho}{n(1-\rho)^2}
  +
  \frac{\rho^n}{1-\rho}
 \right].
 \label{eq:s2-lrv-bound}
\end{equation}
\end{theorem}

\begin{proof}
Stationarity gives the exact finite-\(n\) identity
\[
 D(n)
 =
 \gamma_0(f)
 +
 2\sum_{k=1}^{n-1}
 \left(1-\frac{k}{n}\right)\gamma_k(f).
\]
Absolute summability follows from \eqref{eq:s2-cov-envelope}.  Subtracting the
infinite covariance series yields
\[
 D(n)-\sigma_P^2(f)
 =
 -\frac2n\sum_{k=1}^{n-1}k\gamma_k(f)
 -
 2\sum_{k=n}^{\infty}\gamma_k(f).
\]
Therefore
\[
 |D(n)-\sigma_P^2(f)|
 \leq
 \frac{2C}{n}\sum_{k=1}^{n-1}k\rho^k
 +
 2C\sum_{k=n}^{\infty}\rho^k.
\]
Using
\(\sum_{k\geq1}k\rho^k=\rho/(1-\rho)^2\) and
\(\sum_{k\geq n}\rho^k=\rho^n/(1-\rho)\) proves
\eqref{eq:s2-lrv-bound}.
\end{proof}

\begin{remark}[Nonstationary paths]
\label{rem:s2-nonstationary-horizon}
Theorem~\ref{thm:s2-lrv-envelope} is deliberately restricted to
stationarity.  For \(X_0\sim\nu_a\neq\pi_h\), the finite-horizon objects are
\(D_{a,r}(n)\), \(R_{a,r}(n)\), and \(B_{a,r}(n)\) from
\eqref{eq:s2-dispersion}--\eqref{eq:s2-bias}.  A visually flat dispersion
curve does not remove the initialization term \(nB_{a,r}(n)^2\), and it does
not verify \eqref{eq:s2-cov-envelope}.  Unless usable, independently justified
constants \(C,\rho\) and a stationary initialization are available, Gate 3 is
a falsification screen only.  The decision target must then remain a declared
finite-horizon dispersion, finite-horizon MSE, or bounded finite-horizon loss.
\end{remark}

\subsection{Exact bounded finite-horizon risk replicates}
\label{subsec:stage2-risk-replicates}

We first give exact bounded component-risk replicates.  They do not require
stationarity.

\begin{proposition}[Exact paired finite-horizon dispersion replicate]
\label{prop:s2-paired-dispersion}
Suppose \(m_r\leq f_r\leq M_r\) and
\(\Delta_r=M_r-m_r>0\).  Generate two independent length-\(n\) paths from the
same candidate \(P_a\) and the same initialization law \(\nu_a\), and let
\(\widehat\mu^{(1)}_{a,r,n}\) and
\(\widehat\mu^{(2)}_{a,r,n}\) be their averages.  Then
\begin{equation}
 Z_{a,r}
 =
 \frac n2
 \left(
  \widehat\mu^{(1)}_{a,r,n}
  -
  \widehat\mu^{(2)}_{a,r,n}
 \right)^2
 \label{eq:s2-paired-z}
\end{equation}
satisfies
\begin{equation}
 \mathbb E_{\nu_a}Z_{a,r}=D_{a,r}(n),
 \qquad
 0\leq Z_{a,r}\leq\frac{n\Delta_r^2}{2}.
 \label{eq:s2-paired-range}
\end{equation}
Thus
\[
 Y_{a,r}
 =
 \frac{2Z_{a,r}}{n\Delta_r^2}
 \in[0,1],
 \qquad
 \mathcal J_r(a):=\mathbb EY_{a,r}
 =
 \frac{2D_{a,r}(n)}{n\Delta_r^2}.
\]
\end{proposition}

\begin{proof}
If \(U,V\) are independent and identically distributed, then
\(\mathbb E(U-V)^2=2\operatorname{Var}(U)\).  Apply this with the two path
averages and multiply by \(n/2\).  Each average lies in \([m_r,M_r]\), so its
difference from the other average has magnitude at most \(\Delta_r\), proving
the range.
\end{proof}

\begin{proposition}[Bounded MSE and capped-loss replicates]
\label{prop:s2-bounded-loss}
If \(m_r\leq f_r\leq M_r\) and \(\mu_r\) is known, then
\[
 Y_{a,r}^{\mathrm{MSE}}
 =
 \frac{(\widehat\mu_{a,r,n}-\mu_r)^2}{\Delta_r^2}
 \in[0,1],
 \qquad
 \mathbb E Y_{a,r}^{\mathrm{MSE}}
 =
 \frac{R_{a,r}(n)}{n\Delta_r^2}.
\]
For an unbounded estimand, fix a scientifically declared tolerance
\(\tau_r>0\) before candidate evaluation and define
\[
 \ell_{\tau_r}(e)
 =
 \min\left\{\frac{e^2}{\tau_r^2},1\right\},
 \qquad
 Y_{a,r}^{\mathrm{cap}}
 =
 \ell_{\tau_r}(\widehat\mu_{a,r,n}-\mu_r).
\]
Then \(Y_{a,r}^{\mathrm{cap}}\in[0,1]\) and its expectation is a bounded
finite-horizon, estimand-specific decision loss.  It is not the raw MSE
\(R_{a,r}(n)\).  If a certified reference \(\widetilde\mu_r\) obeys
\(|\widetilde\mu_r-\mu_r|\leq b_r\), then
\begin{equation}
 \left|
  \mathbb E\ell_{\tau_r}
    (\widehat\mu_{a,r,n}-\widetilde\mu_r)
  -
  \mathbb E\ell_{\tau_r}
    (\widehat\mu_{a,r,n}-\mu_r)
 \right|
 \leq
 \frac{2b_r}{\tau_r}.
 \label{eq:s2-reference-loss}
\end{equation}
\end{proposition}

\begin{proof}
The bounded-MSE statement follows because both the path average and its target
mean lie in \([m_r,M_r]\).  The capped loss is bounded by construction.  The
map \(e\mapsto\ell_{\tau_r}(e)\) is globally \(2/\tau_r\)-Lipschitz: on
\(|e|<\tau_r\) its derivative has magnitude at most \(2/\tau_r\), and outside
that interval it is constant.  Equation \eqref{eq:s2-reference-loss} follows
from the Lipschitz bound.
\end{proof}

\subsection{Operational finite-sample selection and certified fallback}
\label{subsec:stage2-selection}
Let
\[
 \mathcal F_0
 =
 \mathcal F_{\mathrm{freeze}}
 \vee\mathcal F_{\mathrm{adq}}
 \vee\mathcal F_{\mathrm{hor}}
\]
contain all information used before risk tuning.  The eligible library
\(\mathcal A_{\mathrm{adq}}\) is \(\mathcal F_0\)-measurable.  Conditional on
\(\mathcal F_0\), suppose the tuning replicates
\(\{Y_{a,r,\ell}:1\leq\ell\leq m_{\mathrm{tun}}\}\) are independent over
\(\ell\), lie in \([0,1]\), and have conditional mean
\(\mathcal J_r(a)\).  Dependence across candidates or components at a fixed
replicate index is allowed.

\begin{theorem}[Operational finite-library selection]
\label{thm:s2-operational-selection}
Let \(K=|\mathcal A|\), let \(R\) be the number of mandatory component risks,
and let \(\delta_{\mathrm{tun}}\in(0,1)\).  For
\[
 \widehat{\mathcal J}_r(a)
 =
 \frac1{m_{\mathrm{tun}}}
 \sum_{\ell=1}^{m_{\mathrm{tun}}}Y_{a,r,\ell},
\]
define
\begin{equation}
 \varepsilon_{\mathrm{tun}}
 =
 \sqrt{
  \frac{\log(2KR/\delta_{\mathrm{tun}})}
       {2m_{\mathrm{tun}}}
 }.
 \label{eq:s2-tuning-epsilon}
\end{equation}
Then, conditionally on \(\mathcal F_0\), with probability at least
\(1-\delta_{\mathrm{tun}}\),
\begin{equation}
 \sup_{a\in\mathcal A_{\mathrm{adq}}}
 \max_{1\leq r\leq R}
 \left|
  \widehat{\mathcal J}_r(a)-\mathcal J_r(a)
 \right|
 \leq\varepsilon_{\mathrm{tun}}.
 \label{eq:s2-uniform-components}
\end{equation}

Let nonnegative weights satisfy \(\sum_r\omega_r=1\), and define
\[
 \mathcal J_{\mathrm w}(a)=\sum_r\omega_r\mathcal J_r(a),
 \qquad
 \mathcal J_{\max}(a)=\max_r\mathcal J_r(a),
\]
with plug-in versions formed from the component estimates.  For either
\(\mathcal J\in\{\mathcal J_{\mathrm w},\mathcal J_{\max}\}\), if
\[
 \widehat{\mathcal J}(\widehat a)
 \leq
 \min_{a\in\mathcal A_{\mathrm{adq}}}
 \widehat{\mathcal J}(a)
 +
 \eta_{\mathrm{opt}},
\]
then on \eqref{eq:s2-uniform-components},
\begin{equation}
 \mathcal J(\widehat a)
 \leq
 \min_{a\in\mathcal A_{\mathrm{adq}}}\mathcal J(a)
 +
 2\varepsilon_{\mathrm{tun}}
 +
 \eta_{\mathrm{opt}}.
 \label{eq:s2-oracle-regret}
\end{equation}
If capped losses use certified references with
\(\kappa_{\mathrm{ref}}=\max_r2b_r/\tau_r\), replace
\(\varepsilon_{\mathrm{tun}}\) in
\eqref{eq:s2-uniform-components}--\eqref{eq:s2-oracle-regret} by
\(\varepsilon_{\mathrm{tun}}+\kappa_{\mathrm{ref}}\).
\end{theorem}

\begin{proof}
For a fixed pair \((a,r)\), Hoeffding's inequality gives
\[
 \mathbb P\left(
  |\widehat{\mathcal J}_r(a)-\mathcal J_r(a)|>t
  \,\middle|\,\mathcal F_0
 \right)
 \leq2e^{-2m_{\mathrm{tun}}t^2}.
\]
Take \(t=\varepsilon_{\mathrm{tun}}\) and apply the union bound over the
at most \(KR\) frozen pairs.  This proves
\eqref{eq:s2-uniform-components}.  On that event, the weighted objective
differs from its plug-in version by at most
\(\sum_r\omega_r\varepsilon_{\mathrm{tun}}
=\varepsilon_{\mathrm{tun}}\), and the maximum objective differs by at most
the same amount.  If \(a^\star\) minimizes the true objective over the eligible
library, then
\[
 \mathcal J(\widehat a)
 \leq
 \widehat{\mathcal J}(\widehat a)+\varepsilon_{\mathrm{tun}}
 \leq
 \widehat{\mathcal J}(a^\star)
 +\eta_{\mathrm{opt}}+\varepsilon_{\mathrm{tun}}
 \leq
 \mathcal J(a^\star)
 +2\varepsilon_{\mathrm{tun}}+\eta_{\mathrm{opt}}.
\]
The reference-error modification follows from
\eqref{eq:s2-reference-loss} and another triangle inequality.
\end{proof}

\begin{remark}[Why this theorem is operational]
\label{rem:s2-operational}
The range \([0,1]\) is analytic and fixed before tuning, so its envelope has
failure probability zero.  No tuning-stream standard deviation is inserted
into the theorem.  The price is conservatism, quantified below.  Raw
finite-horizon dispersion or MSE for an unbounded estimand is not covered by
Theorem~\ref{thm:s2-operational-selection}; it requires either a proved
finite-moment robust construction with an independent scale stream or the
scientifically frozen capped loss in
Proposition~\ref{prop:s2-bounded-loss}.  If neither is approved, the full
portfolio is not operational and the fallback is mandatory.
\end{remark}

After tuning, let
\[
 \mathcal F_{\mathrm{sel}}
 =
 \mathcal F_0\vee\mathcal F_{\mathrm{tun}}
\]
contain the selected candidate and all abort decisions.  Certification data
are conditionally independent of \(\mathcal F_{\mathrm{sel}}\).

\begin{theorem}[Independent certification and fallback]
\label{thm:s2-fallback}
Let \(a_0\) be the frozen fallback and let \(\widehat a\) be
\(\mathcal F_{\mathrm{sel}}\)-measurable.

\emph{Weighted objective.}
For each certification replicate, let
\[
 W_\ell
 =
 \sum_{r=1}^R
 \omega_r
 \{Y_{a_0,r,\ell}-Y_{\widehat a,r,\ell}\}.
\]
Assume the replicates are independent over \(\ell\), conditional on
\(\mathcal F_{\mathrm{sel}}\), with the correct component marginals.  Then
\(W_\ell\in[-1,1]\) and
\[
 \mathbb E(W_\ell\mid\mathcal F_{\mathrm{sel}})
 =
 \Delta_{\widehat a}
 :=
 \mathcal J_{\mathrm w}(a_0)
 -
 \mathcal J_{\mathrm w}(\widehat a).
\]
The lower bound
\begin{equation}
 L_{\mathrm{cert}}
 =
 \overline W
 -
 \sqrt{
  \frac{2\log(1/\delta_{\mathrm{cert}})}
       {m_{\mathrm{cert}}}
 }
 \label{eq:s2-cert-lcb}
\end{equation}
satisfies
\begin{equation}
 \mathbb P\{
  L_{\mathrm{cert}}\leq\Delta_{\widehat a}
  \mid\mathcal F_{\mathrm{sel}}
 \}
 \geq1-\delta_{\mathrm{cert}}.
 \label{eq:s2-cert-coverage}
\end{equation}

\emph{Maximum objective.}
For \(1\leq r,s\leq R\), form
\[
 W_{r,s,\ell}
 =
 Y_{a_0,r,\ell}-Y_{\widehat a,s,\ell}
 \in[-1,1]
\]
and
\[
 L_{r,s}
 =
 \overline W_{r,s}
 -
 \sqrt{
  \frac{2\log(R^2/\delta_{\mathrm{cert}})}
       {m_{\mathrm{cert}}}
 }.
\]
Then
\begin{equation}
 L_{\max}
 =
 \max_r\min_sL_{r,s}
 \label{eq:s2-max-lcb}
\end{equation}
is a simultaneous lower confidence bound for
\(\mathcal J_{\max}(a_0)-\mathcal J_{\max}(\widehat a)\) with conditional
coverage at least \(1-\delta_{\mathrm{cert}}\).

Deploy \(\widehat a\) only if no earlier abort occurred and the applicable
lower bound is strictly positive; otherwise deploy \(a_0\).  Then, for either
objective,
\begin{equation}
 \mathbb P\{
  \mathcal J(\textnormal{deployed})
  >
  \mathcal J(a_0)
 \mid\mathcal F_{\mathrm{sel}}
 \}
 \leq\delta_{\mathrm{cert}}.
 \label{eq:s2-no-harm}
\end{equation}
For capped losses evaluated with an approximate certified reference, subtract
\(2\kappa_{\mathrm{ref}}\) from the lower bounds.
\end{theorem}

\begin{proof}
For the weighted objective, Hoeffding's one-sided inequality for variables in
an interval of length two gives
\[
 \mathbb P\left\{
  \overline W-\Delta_{\widehat a}\leq-t
  \,\middle|\,\mathcal F_{\mathrm{sel}}
 \right\}
 \leq
 \exp\left(-\frac{m_{\mathrm{cert}}t^2}{2}\right).
\]
Substitution of the radius in \eqref{eq:s2-cert-lcb} proves
\eqref{eq:s2-cert-coverage}.

For the maximum objective, the same inequality with error allocation
\(\delta_{\mathrm{cert}}/R^2\) holds simultaneously for all \((r,s)\).
On that event \(L_{r,s}\leq\mathcal J_r(a_0)-\mathcal J_s(\widehat a)\), and
\[
 \mathcal J_{\max}(a_0)-\mathcal J_{\max}(\widehat a)
 =
 \max_r\min_s
 \{\mathcal J_r(a_0)-\mathcal J_s(\widehat a)\}.
\]
Taking \(\max_r\min_s\) preserves the lower-bound direction and proves the
claim for \(L_{\max}\).

On an abort or nonpositive-bound branch the fallback is deployed and cannot
have risk larger than itself.  On a positive-bound branch, harmful deployment
requires failure of the relevant lower-confidence event.  This proves
\eqref{eq:s2-no-harm}.  Under an approximate reference, each candidate's
component risk can move by at most \(\kappa_{\mathrm{ref}}\), so a baseline
minus candidate difference can move by at most
\(2\kappa_{\mathrm{ref}}\).
\end{proof}

\begin{proposition}[Certification power calculation]
\label{prop:s2-cert-power}
For the weighted construction, suppose the true normalized improvement is
\(\Delta_{\widehat a}\geq d>0\).  To make
\(\mathbb P(L_{\mathrm{cert}}>0\mid\mathcal F_{\mathrm{sel}})
\geq1-\beta\), it is sufficient that
\begin{equation}
 m_{\mathrm{cert}}
 \geq
 \frac{2}{d^2}
 \left\{
  \sqrt{\log(1/\delta_{\mathrm{cert}})}
  +
  \sqrt{\log(1/\beta)}
 \right\}^2.
 \label{eq:s2-cert-power-count}
\end{equation}
\end{proposition}

\begin{proof}
With conditional probability at least \(1-\beta\),
\[
 \overline W
 \geq
 \Delta_{\widehat a}
 -
 \sqrt{\frac{2\log(1/\beta)}{m_{\mathrm{cert}}}}.
\]
Combining this inequality with \eqref{eq:s2-cert-lcb}, the lower bound is
positive whenever the two radii sum to less than \(d\).
Solving for \(m_{\mathrm{cert}}\) gives
\eqref{eq:s2-cert-power-count}.
\end{proof}

\subsection{Error budget and executable Tier-A constants}
\label{subsec:stage2-budget}

\begin{proposition}[Complete workflow error allocation]
\label{prop:s2-error-budget}
Let \(\mathcal E_{\mathrm{tar}}\), \(\mathcal E_{\mathrm{adq}}\),
\(\mathcal E_{\mathrm{tun}}\), \(\mathcal E_{\mathrm{env}}\),
\(\mathcal E_{\mathrm{cert}}\), and \(\mathcal E_{\mathrm{rep}}\) be,
respectively, the target/reference-validity, simultaneous Gate-2, uniform
tuning, risk-envelope, independent certification, and held-out-reporting
events.  Suppose their conditional constructions are valid at
failure levels
\[
 \delta_{\mathrm{tar}},\quad
 \delta_{\mathrm{adq}},\quad
 \delta_{\mathrm{tun}},\quad
 \delta_{\mathrm{env}},\quad
 \delta_{\mathrm{cert}},\quad
 \delta_{\mathrm{rep}}.
\]
Sample splitting is used so that each conditional statement is made after all
objects it conditions on have been frozen.  Then
\begin{equation}
 \mathbb P\left(
  \bigcap_{\star\in\{\mathrm{tar,adq,tun,env,cert,rep}\}}
  \mathcal E_\star
 \right)
 \geq
 1-
 \sum_\star\delta_\star.
 \label{eq:s2-total-budget}
\end{equation}
On this intersection, every eligible candidate satisfies the Gate-2
requirement, the selected candidate satisfies
\eqref{eq:s2-oracle-regret}, and a deployed nonfallback candidate improves on
the fallback.  Independence among the five events is not required for
\eqref{eq:s2-total-budget}; the separate streams are required to make the
conditional confidence statements valid.
\end{proposition}

\begin{proof}
Each construction controls its failure probability conditionally on the
sigma-field generated by earlier stages.  Iterated expectation gives the same
unconditional bound.  A union bound proves \eqref{eq:s2-total-budget}.  The
three substantive conclusions are exactly those of
Theorems~\ref{thm:s2-gate2}, \ref{thm:s2-operational-selection}, and
\ref{thm:s2-fallback} on their respective events.
\end{proof}

For the next experiment, freeze the familywise budget
\begin{equation}
\begin{array}{c|cccccc}
\text{stage}
&\text{target}
&\text{adequacy}
&\text{tuning}
&\text{envelope}
&\text{certification}
&\text{held-out report}\\ \hline
\delta
&0.005&0.010&0.010&0&0.015&0.010 .
\end{array}
\label{eq:s2-numeric-budget}
\end{equation}
The envelope allocation is zero because the operational construction uses
analytic \([0,1]\) ranges.  The allocations sum to \(0.05\).

\begin{proposition}[Tier-A numerical instantiation]
\label{prop:s2-tier-a}
Use the three-candidate family \(K=3\).  In the Tier-A mode-discovery rule,
every retained mode has reference mass at least \(0.015\), so disjointness
implies at most
\[
 M_{\mathrm{mode}}
 \leq\left\lfloor\frac1{0.015}\right\rfloor=66
\]
retained modes.  With the five non-mode estimands, \(R\leq71\); with two
non-mode nonsmooth events, the number of mandatory nonsmooth estimands is at
most \(68\), and hence
\[
 M_{\mathrm{adq}}\leq3\times68=204.
\]
Under equal Bonferroni allocation, \(z^{\max}=0.10\), and
\(\delta_{\mathrm{adq}}=0.010\), zero observed zero-crossing paths require
\begin{equation}
 L_{\mathrm{adq}}
 \geq
 \left\lceil
 \frac{\log(0.010/204)}{\log(0.9)}
 \right\rceil
 =
 95.
 \label{eq:s2-tier-a-adq}
\end{equation}
For the fixed seven-estimand Tier-B portfolio with four nonsmooth components,
the corresponding count is
\[
 \left\lceil
 \frac{\log(0.010/12)}{\log(0.9)}
 \right\rceil
 =
 68.
\]

For \(K=3\), \(R=71\), and \(\delta_{\mathrm{tun}}=0.010\), requiring
\(\varepsilon_{\mathrm{tun}}\leq0.10\) in
\eqref{eq:s2-tuning-epsilon} requires
\begin{equation}
 m_{\mathrm{tun}}
 \geq
 \left\lceil
  \frac{\log(2\cdot3\cdot71/0.010)}
       {2(0.10)^2}
 \right\rceil
 =
 533
 \quad\text{replicates per candidate}.
 \label{eq:s2-tier-a-tuning}
\end{equation}
For the fixed seven-component portfolio the corresponding count is \(418\).
The legacy value \(m_{\mathrm{tun}}=2\) gives
\(\varepsilon_{\mathrm{tun}}=1.63245\ldots\) in the worst-case Tier-A
portfolio and is vacuous.

Finally, with \(\delta_{\mathrm{cert}}=0.015\), desired power \(0.80\), and a
scientifically material normalized improvement \(d=0.20\),
\eqref{eq:s2-cert-power-count} gives
\begin{equation}
 m_{\mathrm{cert}}\geq551
 \label{eq:s2-tier-a-cert}
\end{equation}
paired certification replicates for the weighted objective.  These values are
finite and executable, but intentionally expose the cost of an assumption-free
bounded-range guarantee.
\end{proposition}

\begin{proof}
The mode-count bound follows because \(m\) disjoint sets of probability at
least \(0.015\) have total probability at least \(0.015m\leq1\).  The remaining
values follow by direct substitution into
\eqref{eq:s2-required-adq}, \eqref{eq:s2-tuning-epsilon}, and
\eqref{eq:s2-cert-power-count}; the accompanying verification script performs
the calculations without rounding intermediate values.
\end{proof}

\begin{remark}[Remaining Tier-A target-validity hold]
\label{rem:s2-tier-a-hold}
The numerical counts in Proposition~\ref{prop:s2-tier-a} instantiate the
selection theorem's library size, multiplicity, concentration radii, and power.
They do not fabricate missing componentwise quadrature errors.  Before a
Tier-A capped-loss experiment is run, the tolerance vector \(\tau_r\) and the
reference radii \(b_r\) must be frozen and must satisfy the target budget
\(\delta_{\mathrm{tar}}=0.005\).  The natural provisional scales are
\(\tau_\alpha=0.5\), the prior standard deviation of log-intensity, and
\(\tau_{p_\star}=0.195\), the frozen regulatory concentration threshold.
These become operational only after PI approval and an independent
componentwise reference-convergence table.  Until then, the full-portfolio
selection procedure remains stopped.
\end{remark}
